\subsection{The excess-one curve calculation and the ordinary-double sector}

\begin{statement}{P4.40}{Proved}
In the scope of Part~4.2 where every fibre support is
irreducible or normal crossing, no irreducible support with
positive-genus ruled resolution has generic transverse type
$A_1$ on every component of its critical curve.
\end{statement}
\begin{proof}
We first give the exact local curve calculation needed in
the proof. Let $A$ be a reduced complex curve local ring with
$r$ branches, let $B=\bigoplus_{i=1}^r\mathbb C\{t_i\}$
be its normalization, and let $R\subset B$ be the ring whose
constant terms agree. Then
\[
 \operatorname{length}(R/A)=\delta(A)-(r-1).
\]
Suppose this length is one. The simple $A$-module $R/A$
is annihilated by $\mathfrak m_A$, so
$\mathfrak m_A R\subset A$. Choose
$v\in\mathfrak m_R$ spanning $R/A$ and write
$v^2=a+cv$, where $a\in\mathfrak m_A$ and $c\in\mathbb C$.
If $c\ne0$, the convergent square-root series solves
$b^2-cb-a=0$ with $b\in\mathfrak m_A$. Then
$(v-b)(v+b-c)=0$ and the second factor is a unit, forcing
$v=b$, a contradiction. Thus $v^2\in A$ and
\[
                         \mathfrak m_R^2\subset A.
\]
Consequently $A$ is exactly the inverse image of a hyperplane
\[
                         \sum_i a_i f_i'(0)=0
\]
in $\mathfrak m_R/\mathfrak m_R^2=\mathbb C^r$.
Rescaling the parameters with nonzero $a_i$ gives actual
analytic ring isomorphisms with the following models.

With one nonzero coefficient, that branch has ring
$\mathbb C\{t^2,t^3\}$ and the others are independent smooth
branches joined at their constant terms. The embedding
dimension is $r+1$: for $r=1$ this is a plane cusp, for
$r=2$ a cusp and transverse line, and for $r=3$ it does
not embed in a smooth threefold. With two nonzero
coefficients, the pair has equal first derivatives. Its
ring is generated by $(t,t)$ and $(t^2,-t^2)$, giving the
plane tacnode $y^2=x^4$. The other $r-2$ branches are
independent transverse lines. For $r=3$ and three nonzero
coefficients, the first jets are the restrictions of linear
forms to three distinct lines in a plane. Their products
span all three second-jet coordinates, so the ring is
exactly the ordinary plane triple point. Equivalently,
the latter ring has precisely the same constant and
first-derivative conditions and colength three in $B$.
These cases exhaust $r\le3$. The analytic ring
isomorphisms lift to ambient coordinate changes: choose
minimal algebra generators, lift their independent linear
parts, and apply the inverse-function theorem. In embedding
dimension two we first straighten the smooth containing
surface. Every realizable model is therefore one of the
local classes excluded in P4.38 for ordinary-double
transverse type.

Now suppose the stated ruled support exists. Its horizontal
pole section has coefficient two. It cannot map
nonconstantly to the conductor, since an ordinary double
point has conductor exponent one on each normalization
branch. The section is contracted. P4.33 gives
$\chi(\mathcal O_{C_{\rm sch}})=-1$.
The conductor image is Cohen--Macaulay and generically
reduced in type $A_1$. It is therefore reduced, so
$C_{\rm sch}=C$ and $\chi(\mathcal O_C)=-1$.
All normalized components of $C$ are rational.

The already proved group and Euler calculations leave
exactly two normalized rational components and one
point-identification cycle. Indeed the rational-tree case
has Euler number at least three by P4R.4; the
one-component cycle is excluded by P4.35; and the global
nonnegative Euler sum is at most two. The normalization
sequence and the graph formula consequently give
\[
 \sum_p\delta_p=3,\qquad
 \sum_p(r_p-1)=2,
 \qquad
 \sum_p\bigl(\delta_p-(r_p-1)\bigr)=1.
\]
Every summand in the last sum is nonnegative. There is
exactly one excess-one point and every point has at most
three local branches. The local calculation classifies that
point by an actual analytic isomorphism, and P4.38 excludes
each possible model. This is the required contradiction.
\end{proof}

The first-jet cases, the second-jet product determinant, and
the global excess computation are verified in
\texttt{code/part4/curve\_excess.py}; its saved output reports
PASS. The proof retains the ordinary-double boundary
condition and makes no substitution of transverse type
$A_2$ for $A_1$.
